\documentclass[12pt,a4paper]{article}
\usepackage{amsmath,amssymb}
\usepackage{enumitem}
\usepackage{hyperref}

\title{\bf AI Prompts for Bridging Gap~3\\
\large From the Dirac Operator with Prime Point Interactions\\
to the Weil Explicit Formula}
\author{}
\date{}

\begin{document}
\maketitle

The following prompts are designed to be given sequentially to a suitably capable AI (or used as a research guide) in order to attack Gap~3: proving that the spectral action of the self‑adjoint Dirac operator \(D_\infty\) equals the Weil explicit formula.

\begin{enumerate}[leftmargin=*,label=\textbf{Prompt \arabic*:},ref=Prompt \arabic*]

\item \textbf{Exact determinant for finite \(\lambda\).}\\
Derive an explicit expression for the regularised spectral determinant \(\Delta_\lambda(z)\) of the Dirac operator \(D_\lambda\) on a circle of length \(L = 2\log\lambda\) with point interactions at \(x_p = \log p \bmod L\) (for all primes \(p \le \lambda^2\)), where the matching condition at each point is \(\psi(x_p^+) = e^{i\alpha_p}\psi(x_p^-)\) with \(\alpha_p = (\log p)/\sqrt{p}\).
Use the Krein resolvent formula and transfer matrix methods to obtain a finite product representation involving the coupling constants and the free Dirac propagator.
Express the result in a form that makes the connection to the truncated Weil quadratic form of Connes--Consani--Moscovici (CCM) apparent.

\item \textbf{Relating the finite determinant to the CCM secular equation.}\\
The CCM spectral triple involves a finite‑dimensional matrix \(Q\) (the Weil quadratic form) on \(E_N(\lambda) = \operatorname{span}\{e^{i\mu_k x}\}_{|k|\le N}\).
Prove that the eigenvalues of \(D_\lambda\) coincide, up to a controlled error vanishing as \(\lambda, N\to\infty\) with \(L\asymp\log N\), with the roots of \(\det(Q - zI)\).
Show that the regularised determinant \(\Delta_\lambda(z)\) is equal to \(\det(Q - zI)\) multiplied by an explicit Gamma factor arising from the archimedean contribution.
This equivalence should be established by comparing the boundary‑value problem for \(D_\lambda\) with the CCM eigenvalue equation \(Q\xi = \epsilon\xi\) and using the Fourier basis.

\item \textbf{Asymptotics of the determinant as \(\lambda\to\infty\).}\\
Assume Conjecture~1 of CCM (the ground‑state eigenvector coefficients are given by a prime sum with power‑of‑log error).
Use this to prove that, after multiplying by a suitable product of Gamma functions (the local archimedean factor), the finite determinants \(\Delta_\lambda(z)\) converge uniformly on compact subsets of \(\mathbb{C}\) to the completed Riemann \(\Xi\)-function:
\[
\lim_{\lambda\to\infty} \Delta_\lambda(z) \cdot \Gamma_{\text{arch}}(z) = \Xi(z).
\]
The proof should combine the explicit prime‑sum expression for the eigenvector coefficients with the infinite product representation of \(\Xi(z)\) (Hadamard product over zeros) and the Weil explicit formula to control the difference.

\item \textbf{Contour integration and spectral action.}\\
For a self‑adjoint operator \(A\) with purely discrete spectrum and known regularised determinant \(\det_{\text{reg}}(A - z)\), derive the formula
\[
\operatorname{Tr}(h(A)) - \operatorname{Tr}(h(A_0)) = \frac{1}{2\pi i} \oint_{\mathcal{C}} h(z) \frac{d}{dz} \log \frac{\det_{\text{reg}}(A - z)}{\det_{\text{reg}}(A_0 - z)} \, dz,
\]
where \(A_0\) is a reference operator and \(h\) is analytic in a neighbourhood of the spectrum.
Apply this to \(A = D_\infty\) and \(A_0 = D_0\) (the free Dirac operator), assuming the convergence of the determinants established in Prompt~3.
Evaluate the contour integral to obtain a sum over the eigenvalues of \(D_\infty\) in terms of the logarithmic derivative of \(\Xi(z)\).

\item \textbf{Final identification with the Weil explicit formula.}\\
Starting from the identity
\[
\operatorname{Tr}(h(D_\infty)) = \frac{1}{2\pi i} \oint_{\mathcal{C}} h(z) \frac{\Xi'(z)}{\Xi(z)} \, dz + \text{(trivial zero contributions)},
\]
use the classical Weil explicit formula (which expresses the sum over primes in terms of an integral of \(h\) against \(\Xi'/\Xi\) plus archimedean terms) to rewrite the right‑hand side as
\[
\sum_{p} \frac{\log p}{\sqrt{p}} h(\log p) + \text{(archimedean and trivial zero terms)}.
\]
Show that the contour integral exactly recovers the zero‑sum \(\sum_{\rho} h(\rho)\) and that the remaining terms match those in the explicit formula.
Conclude that the trace formula (Gap~3) holds for all admissible test functions.

\end{enumerate}

These five prompts, when answered in sequence, collectively bridge Gap~3.
The core difficulty lies in Prompt~3 (the asymptotics), which is equivalent to proving the Riemann Hypothesis; the other prompts are technical but manageable.

\end{document}